% Generated draft: requires mathematical and source review.
\documentclass[11pt]{article}
\usepackage[margin=25mm]{geometry}
\usepackage{fontspec}
\setmainfont{Latin Modern Roman}
\usepackage{amsmath,amssymb,mathtools,unicode-math}
\setmathfont{Latin Modern Math}
\usepackage{enumitem,needspace,xurl,hyperref}
\setlength{\parindent}{0pt}
\setlength{\parskip}{5pt}
\setlength{\emergencystretch}{3em}
\newcommand{\N}{\mathbb{N}}
\newcommand{\Z}{\mathbb{Z}}
\newcommand{\Q}{\mathbb{Q}}
\newcommand{\R}{\mathbb{R}}
\newcommand{\C}{\mathbb{C}}
\newcommand{\F}{\mathbb{F}}
\newcommand{\A}{\mathbb{A}}
\newcommand{\PP}{\mathbb{P}}
\newcommand{\E}{\mathbb{E}}
\newcommand{\eps}{\varepsilon}
\newcommand{\dd}{\,\mathrm{d}}
\begin{document}
\section*{UnsolvedMath ID: 30004218}
\textbf{Source problem code:} catalogue entry
\subsection{Definitions and mathematical statement}
For a map \(f:X_f\to Y_f\) and a point \(y\in Y_f\), its fiber is
\[
F_{f,y}=f^{-1}(\{y\})=\{x\in X_f:f(x)=y\},
\]
with the subspace topology. For an integer \(d\geq 0\), write
\[
\overline B^{d}=\{u\in\R^{d}:\|u\|\leq 1\},
\]
where \(\overline B^0\) is a single point. A CW complex is regular if every characteristic map from a closed ball onto the closure of a cell is a homeomorphism. Thus, asking whether a specified cell decomposition is a regular CW decomposition concerns that particular decomposition, including the CW topology and attaching maps; it requires more than the existence of some regular CW structure on the underlying space.

The supplied record does not specify the maps, their domains and targets, or the natural cell decompositions. Let \(\mathcal M\) provisionally denote the intended family of maps, and let \(\mathcal D_{f,y}\) denote the intended decomposition of \(F_{f,y}\). These symbols are placeholders for missing source definitions, not a newly chosen family or decomposition.

The intended question is whether, for every intended map \(f\in\mathcal M\) and every intended fiber \(F_{f,y}\), the decomposition \(\mathcal D_{f,y}\) gives \(F_{f,y}\) the structure of a regular CW complex and there is an integer \(d=d(f,y)\geq 0\) such that
\[
F_{f,y}\cong\overline B^{d},
\]
where \(\cong\) denotes homeomorphism. The dimension may depend on the fiber; the record supplies no dimension formula.

\textbf{Unresolved scope.} The phrase ``maps to totally nonnegative spaces'' does not identify a mathematical family of maps or define total nonnegativity in the intended setting. The record also does not identify the cells or their attaching maps. Moreover, it does not state whether ``every fiber'' ranges over all target points or only points in the image. This matters because an empty fiber is not homeomorphic to any \(\overline B^d\) with \(d\geq0\). Consequently, the displayed conclusion is a statement schema requiring the missing source definitions and quantifier range before it becomes a fully specified conjecture.
\subsection{Short English statement}
For each fiber of the intended maps to totally nonnegative spaces, the question asks whether its specified natural cells form a regular CW complex whose underlying space is a closed ball. The supplied record leaves the maps, decompositions, and treatment of empty fibers unspecified, so the precise scope remains unresolved.
\subsection{Sources}
\begin{itemize}
\item Local UnsolvedMath registry: supplied complete record for permanent ID \(30004218\), titled ``Regular Ball Decompositions of Totally Nonnegative Fibers,'' in the Oberwolfach Reports Open Problems collection. Exact source problem code: \texttt{catalogue entry}. The record supplies proposed year \(2019\), but names no proposer.
\item Supplied source citation: \emph{Geometric, Algebraic, and Topological Combinatorics} (2020), Oberwolfach Reports, pp.~2395--2472, DOI \texttt{10.4171/owr/2019/39}. \url{https://doi.org/10.4171/owr/2019/39}. The background describes Oberwolfach Reports as documenting mathematical workshops and their open-problem discussions.
\item Supplied literature source: Davis, Hersh, and Miller, ``fibers of maps'' (the label supplied by the registry). \url{https://arxiv.org/abs/1903.01420}. The record does not provide a full bibliographic citation or explicitly attribute the conjecture to these authors.
\item Supplied literature source: ``Oberwolfach Report 39, fiber conjecture.'' \url{https://doi.org/10.4171/owr/2019/39}.
\item Literature assessment dated \texttt{2026-08-22}: the supplied curation labels the problem open and reports that the fibers have cell stratifications with contractible interior dual block posets, while the stronger regular-CW and closed-ball assertion remains conjectural. This is supplied dated curation derived from the OWR material, not an independently verified assessment; its status and cited claims require independent verification.
\item The registry labels the statement \texttt{corrected\_verified} and reports that the official Oberwolfach report was checked directly, with definitions restored and source-dependent wording removed. These are supplied provenance claims, not verification performed for this draft; the complete record reproduced here still omits the definitions needed to identify the maps and cells.
\end{itemize}
\end{document}
